% ECONOMICS_PROBLEM: {"id": "O33-23", "title": "Selection and attraction of knowledge-growth traveling waves", "jel_code": "O33", "source_id": "OP-0212"}
% !TeX program = xelatex
\documentclass[11pt,a4paper]{article}
\usepackage[margin=23mm]{geometry}
\usepackage{fontspec}
\setmainfont{Latin Modern Roman}
\usepackage{amsmath,amssymb,mathtools}
\usepackage{xcolor,enumitem,booktabs,longtable,array,xurl,hyperref}
\definecolor{muted}{HTML}{53636C}
\hypersetup{colorlinks=true,urlcolor=blue,linkcolor=blue}
\setlength{\parindent}{0pt}
\setlength{\parskip}{5pt}
\newcommand{\R}{\mathbb R}
\newcommand{\E}{\mathbb E}
\newcommand{\N}{\mathbb N}
\newcommand{\ind}{\mathbf 1}
\newcommand{\Var}{\operatorname{Var}}
\newcommand{\Cov}{\operatorname{Cov}}
\newcommand{\supp}{\operatorname{supp}}
\DeclareMathOperator*{\argmin}{arg\,min}
\DeclareMathOperator*{\argmax}{arg\,max}
\newcommand{\entrymeta}[3]{{\sffamily\small\color{muted}\textbf{\texttt{#1}}\quad|\quad #2\par #3\par}}
\newcommand{\Problemref}[1]{\texttt{#1}}
\newcommand{\JELref}[2]{\texttt{#2} (JEL \texttt{#1})}
\begin{document}
\section*{Selection and attraction of knowledge-growth traveling waves}
\textbf{Problem ID:} \texttt{O33-23}\quad \textbf{JEL:} \texttt{O33}\par
% BEGIN ECONOMICS STATEMENT
\label{problem:OP-0212}
% GAME_THEORY_PROBLEM: {"local_id": "GT-082", "original_number": 82, "kind": "R", "entry_kind": "statement_only", "published": false, "review_required": true, "category": "game_theory"}
\label{game:GT-082}
{\sffamily\small\color{muted}
{\small\textbf{GT-082} | Mean-field games | R: precise asymptotic instantiation of a published open question\par}
\par}
\subsection*{Definitions and mathematical statement}
Fix $\nu,\rho>0$ and a continuous increasing concave learning function $\alpha:[0,1]\to[0,\infty)$ with $\alpha(0)=0$. Let $f(t,x)$ be the probability density of log productivity $x\in\mathbb R$. In the diffusive Lucas--Moll model the value $v$, density $f$, and optimal search fraction $s_*\in[0,1]$ solve
\[
 \begin{split}
 -v_t-\nu v_{xx}+\rho v
 &=\max_{s\in[0,1]}\left\{(1-s)e^x+
 \alpha(s)\int_x^\infty[v(t,y)-v(t,x)]f(t,y)dy\right\},\\
 f_t-\nu f_{xx}
 &=f(t,x)\int_{-\infty}^x\alpha(s_*(t,y))f(t,y)dy
   -\alpha(s_*(t,x))f(t,x)\int_x^\infty f(t,y)dy,
 \end{split}
\]
Seek $C^{1,2}$ functions $v,f$ with the displayed integrals finite. Here $s_*(t,x)$ is a measurable maximizing selector, $f_0$ is a positive $C^2$ probability density with $\int e^xf_0(x)dx<\infty$, $f(0)=f_0$, $f\geq0$, and $\int f(t,x)dx=1$. The infinite-horizon boundary condition for $v$ is its value representation
\[
 v(t,x)=\sup_s\mathbb E_{t,x}\int_t^\infty
       e^{-\rho(\tau-t)}(1-s_\tau)e^{X_\tau}d\tau.
\]
Against fixed $f$, $X$ has diffusion $\sqrt{2\nu}\,dB$ and meeting rate $\alpha(s)$; at a meeting with $Y\sim f(\tau,\cdot)$ it becomes $\max(X,Y)$. Controls are progressively measurable search fractions, and the displayed value must be finite. This defines the backward condition instead of imposing a freely chosen $v(0)$.

A balanced-growth traveling wave has
\[
 f(t,x)=\varphi(x-ct),\qquad v(t,x)=e^{ct}V(x-ct),\qquad
 s_*(t,x)=S(x-ct),\quad \int\varphi=1.
\]
Its profile solves
\[
 \begin{split}
 -\nu V''+cV'+(\rho-c)V
 &=\max_{s\in[0,1]}\left\{(1-s)e^z+
 \alpha(s)\int_z^\infty[V(y)-V(z)]\varphi(y)dy\right\},\\
 -c\varphi'-\nu\varphi''
 &=\varphi(z)\int_{-\infty}^z\alpha(S(y))\varphi(y)dy
  -\alpha(S(z))\varphi(z)\int_z^\infty\varphi(y)dy.
 \end{split}
\]
For a concrete nonempty wave regime with a sufficiently large discount, take $\alpha(s)=a\sqrt{s}$, $a>\nu$, and $\rho>a+\nu$. These imply $\rho\geq2\sqrt{\nu a}$, matching the primary paper's $\nu=\kappa^2$ convention. This is a declared subfamily, not the largest known existence regime.

Normalize profiles by $\int_{-\infty}^0\varphi=1/2$. For a solution with a unique median define $b(t)$ by $\int_{-\infty}^{b(t)}f(t,x)dx=1/2$. Characterize the initial densities $f_0$ and solution selections for which some normalized wave obeys
\[
 \frac{b(t)}t\longrightarrow c,\qquad
 \|f(t,b(t)+\cdot)-\varphi\|_{L^1(\mathbb R)}\longrightarrow0,\qquad
 e^{-b(t)}v(t,b(t)+\cdot)\longrightarrow V
 \ \text{in }C^1_{\mathrm{loc}}(\mathbb R).
\]
Determine how the speed and profile depend on $f_0$, including any tail assumptions. This attraction-basin formulation is editorial; arbitrary probability densities are not asserted to satisfy it.
\subsection*{Short English statement}
Which knowledge distributions converge, after translation and value rescaling, to a balanced-growth wave, and which wave is selected?
\subsection*{Variants and scope boundaries}
Wave existence is already known in specified regimes. Convergence in a moving quantile frame allows sublinear shifts, including possible logarithmic shifts; convergence in the rigid frame $x-ct$ is stronger. Finite-horizon terminal values and infinite-horizon value representations are different boundary conditions.
\subsection*{Source relationship and status}
The chapter retains time-dependent convergence as open. The model, moving centering, and norms are explicit specifications; they do not certify the full parameterized claim as open in 2026.
\subsection*{References}
\begingroup\footnotesize\raggedright
{\footnotesize\raggedright
\textbf{[S1]} Cardaliaguet and Delarue. \emph{Selected topics in mean field games}, ICM2022, vol.~5; DOI: 10.4171/ICM2022/17. \textit{Locator:} Section~3.3, ``Traveling waves,'' pp.~3680--3681.\par
\url{https://ems.press/content/book-chapter-files/33265}\par\smallskip
\textbf{[S2]} Alessio Porretta and Luca Rossi (2020). \emph{Traveling waves for a nonlocal KPP equation and mean-field game models of knowledge diffusion}. arXiv:2010.10828v1.\par
\textit{Locators:} system~(2), p.~4; traveling-wave scaling~(3), p.~4; Section~2, Theorem~2.2, critical-wave existence.\par
\url{https://arxiv.org/pdf/2010.10828v1}\par}
\par\endgroup

\subsection*{Common conventions: Source mathematical conventions}

\textbf{Finite games.} For a finite set $A$, $\Delta(A)$ is its probability
simplex. Players choose independent mixed strategies in Nash equilibrium;
correlation is allowed only when explicitly part of the solution concept.
In a game with payoff functions $u_i$ and strategy profile $x$, an additive
$\varepsilon$ Nash equilibrium satisfies
\[
 u_i(x)\geq u_i(x_i',x_{-i})-\varepsilon
 \quad\text{for every player $i$ and every admissible deviation $x_i'$}.
\]
For numerical approximation, payoffs are normalized as locally specified.
An $\varepsilon$ payoff residual is distinct from distance at most
$\varepsilon$ to the set of exact equilibria. Point convergence, convergence
of empirical play, and convergence in distance to an equilibrium set are
also distinct.

\textbf{Computational representations.} Unless stated otherwise, a complexity
question takes rational data with binary encoding; polynomial time is measured
in total bit length. A fixed-accuracy polynomial algorithm is not necessarily
polynomial in $1/\varepsilon$, and neither is necessarily polynomial in
$\log(1/\varepsilon)$. Succinct games and value, demand, or sampling oracles
must declare their interfaces. Exact real solutions require an explicit
representation; an unspecified list of arbitrary real numbers is not a
finite computational output. A claimed algorithmic barrier is meaningful
only for its specified model and reductions.

\textbf{Dynamic games.} Histories, information, observation, and allowable
randomization are part of the model. A stationary strategy depends only on
the current state; a positional strategy in a deterministic graph game
chooses one action at each controlled vertex. Nash equilibrium at an initial
state and equilibrium after every history are different requirements.
An expectation of a pathwise $\liminf$ payoff, a $\liminf$ of expected averages,
a limiting expected average, and a uniform finite-horizon equilibrium are
different criteria. None is silently substituted for another.

\textbf{Allocation and mechanisms.} A complete allocation assigns every item
exactly once unless otherwise stated. Zero-valued goods, negative values,
capacity constraints, feasibility, lotteries, and copies of goods affect
fairness definitions and are specified locally. Dominant-strategy incentive
compatibility, Bayesian incentive compatibility, universal truthfulness,
and truthfulness in expectation impose different inequalities. Whenever
an objective ratio has zero denominator, its convention is stated or those
instances are excluded.

\textbf{Infinite-dimensional models.} Expectations, integrals, stochastic
equations, and optimization objectives require their local regularity,
integrability, filtrations, and admissible domains. A backward--forward
mean-field system is a boundary-value problem; it is not an initial-value
semiflow without an additional construction. Long-time, large-population,
and vanishing-noise limits require an order of limits and a convergence
topology. If a minimum need not be attained, the objective is its infimum
or supremum and the approximation requirement is stated separately.
% END ECONOMICS STATEMENT
\end{document}
